% Generated draft: requires mathematical and source review.
\documentclass[11pt]{article}
\usepackage[margin=25mm]{geometry}
\usepackage{fontspec}
\setmainfont{Latin Modern Roman}
\usepackage{amsmath,amssymb,mathtools,unicode-math}
\setmathfont{Latin Modern Math}
\usepackage{enumitem,needspace,xurl,hyperref}
\setlength{\parindent}{0pt}
\setlength{\parskip}{5pt}
\setlength{\emergencystretch}{3em}
\newcommand{\N}{\mathbb{N}}
\newcommand{\Z}{\mathbb{Z}}
\newcommand{\Q}{\mathbb{Q}}
\newcommand{\R}{\mathbb{R}}
\newcommand{\C}{\mathbb{C}}
\newcommand{\F}{\mathbb{F}}
\newcommand{\A}{\mathbb{A}}
\newcommand{\PP}{\mathbb{P}}
\newcommand{\E}{\mathbb{E}}
\newcommand{\eps}{\varepsilon}
\newcommand{\dd}{\,\mathrm{d}}
\begin{document}
\section*{UnsolvedMath ID: 30002404}
\textbf{Source problem code:} catalogue entry
\subsection{Definitions and mathematical statement}
The supplied question concerns a group \(\Gamma\) acting on a compact space \(X\), with the action described as \emph{asymptotically proper}. The record supplies no definition of this property. It also does not specify the topology of \(\Gamma\), the required properties of \(X\) beyond compactness, the coefficient ring or category, or the degrees of algebraic \(K\)-theory. Consequently, the action hypothesis and the intended assembly map cannot be uniquely formalized from the supplied evidence.

To describe the requested conclusion, consider the following conventional formulation, without asserting that it is the formulation intended by the source. Let \(\Gamma\) be a discrete group, let \(R\) be an associative unital ring, and let \(n\in\Z\). A subgroup is virtually cyclic if it contains a cyclic subgroup of finite index. Write \(\mathcal{VC}(\Gamma)\) for the family of virtually cyclic subgroups. A classifying \(\Gamma\)-CW complex \(E_{\mathcal{VC}}\Gamma\) satisfies
\[
(E_{\mathcal{VC}}\Gamma)^H\simeq *\quad\text{if }H\in\mathcal{VC}(\Gamma),
\qquad
(E_{\mathcal{VC}}\Gamma)^H=\varnothing\quad\text{otherwise}.
\]
Let \(\mathbf K_R\) be the orbit-category spectrum whose value at \(\Gamma/H\) is the nonconnective algebraic \(K\)-theory spectrum of the group ring \(R[H]\). The corresponding assembly map is induced by \(E_{\mathcal{VC}}\Gamma\to *\):
\[
A_n^{\Gamma,R}\colon
H_n^\Gamma(E_{\mathcal{VC}}\Gamma;\mathbf K_R)
\longrightarrow K_n(R[\Gamma]),
\]
where \(H_n^\Gamma\) denotes equivariant homology with these spectrum coefficients. Split injectivity means that there exists a group homomorphism
\[
s_n^{\Gamma,R}\colon K_n(R[\Gamma])\longrightarrow
H_n^\Gamma(E_{\mathcal{VC}}\Gamma;\mathbf K_R)
\quad\text{such that}\quad
s_n^{\Gamma,R}\circ A_n^{\Gamma,R}=\operatorname{id}.
\]

\textbf{Requested statement, subject to unresolved source qualifications.} For an action \(\Gamma\curvearrowright X\) satisfying the source's intended meaning of asymptotic properness, determine whether controlled \(K\)-theory proves split injectivity of the intended algebraic \(K\)-theory Farrell--Jones assembly map. The source also asks whether the Farrell--Jones conjecture follows. In the conventional version displayed above, the full assembly conclusion would require \(A_n^{\Gamma,R}\) to be an isomorphism in every intended degree and for every intended coefficient ring. Split injectivity alone does not establish surjectivity. The missing action definition, coefficient scope, degree quantifiers, and any additional argument or hypotheses needed for the full conjecture remain unresolved; the displayed conventional version does not supply those missing choices.
\subsection{Short English statement}
The supplied question asks whether an asymptotically proper action on a compact space enables controlled \(K\)-theory to prove split injectivity of a Farrell--Jones assembly map, and whether the full conjecture follows. The action condition and assembly parameters are unspecified. Split injectivity alone does not give the isomorphism required by the full conjecture.
\subsection{Sources}
\begin{itemize}
\item Local UnsolvedMath registry: supplied complete record with permanent ID \(30002404\), exact source problem code \texttt{catalogue entry}, title \emph{Controlled K-Theory for Farrell--Jones Assembly Maps}, in the Oberwolfach Reports Open Problems collection. The proposed year is recorded as 2013; no individual proposer is supplied.
\item Supplied source citation: \emph{Noncommutative Geometry (2014)}, Oberwolfach Reports, DOI \texttt{10.4171/owr/2013/45}, pp.~2553--2629. Supplied source link: \url{https://doi.org/10.4171/owr/2013/45}. The record describes Oberwolfach Reports as documenting workshops and their open-problem discussions.
\item Supplied literature source, labelled ``Oberwolfach Report'': \url{https://ems.press/content/serial-article-files/46476}. The supplied literature assessment, checked on \textbf{2026-08-21}, labels the problem open and states that the report's Farrell--Jones results were only outlined and were still being written and checked, and that no authoritative completed proof of the extracted claim was located. This is supplied dated curation, not an independently verified literature assessment.
\item The registry labels the statement \texttt{corrected\_verified} and says that it was checked against the official report to restore hypotheses and notation while preserving the open questions. These are supplied verification claims; this draft does not independently verify them or the supplied open status.
\end{itemize}
\end{document}
